\chapter{REDES NEURAIS PROFUNDAS (\textit{DNN})}

Este capítulo apresenta a base teórica das redes neurais profundas utilizadas nos experimentos. São discutidos os fundamentos do aprendizado profundo, incluindo a arquitetura do Perceptron Multicamadas, as funções de ativação e o algoritmo de retropropagação, a arquitetura de \textit{Autoencoder} aplicada à detecção de anomalias não supervisionada, e as particularidades da aplicação de DNNs ao problema de classificação de tráfego de rede.

\section{Fundamentos de \textit{Deep Learning}}

O aprendizado profundo (\textit{Deep Learning} - DL) é uma área do aprendizado de máquina que utiliza redes neurais com múltiplas camadas para extrair representações hierárquicas a partir dos dados brutos \cite{bochie2020sbrc}. Enquanto nos algoritmos tradicionais de ML a engenharia de características (\textit{feature engineering}) depende do conhecimento prévio do especialista, o aprendizado profundo é capaz de descobrir automaticamente as representações mais discriminativas ao longo das camadas da rede. Esse aprendizado enquadra-se na definição de Mitchell \cite{mitchell1997}: o modelo melhora seu desempenho em uma tarefa à medida que acumula experiência a partir dos dados de treinamento.

A unidade fundamental de uma rede neural artificial é o neurônio artificial, inspirado no neurônio biológico. Cada neurônio recebe um vetor de entradas, aplica uma transformação linear ponderada e passa o resultado por uma função de ativação não linear. Formalmente, a saída do neurônio $i$ na camada $l$ é dada por:
\[
  h_i = \varphi\!\left(W_i^\top h_{i-1} + b_i\right)
\]
onde $W_i$ é o vetor de pesos, $h_{i-1}$ é o vetor de ativações da camada anterior, $b_i$ é o viés e $\varphi$ é a função de ativação \cite{bochie2020sbrc}. Funções de ativação clássicas incluem a sigmoide e a tangente hiperbólica; a \textit{Rectified Linear Unit} (ReLU), definida como $\varphi(x) = \max(0, x)$, tornou-se a escolha padrão em redes profundas por atenuar o problema do gradiente que se dissipa nas camadas mais profundas.

O Perceptron Multicamadas (\textit{Multilayer Perceptron} - MLP) é a arquitetura de referência utilizada neste trabalho. É composto por uma camada de entrada, uma ou mais camadas ocultas e uma camada de saída. O treinamento é realizado pelo algoritmo de retropropagação (\textit{backpropagation}), que calcula o gradiente da função de perda em relação a cada parâmetro da rede por meio da regra da cadeia e atualiza os pesos iterativamente por descida de gradiente estocástica \cite{bochie2020sbrc}.

O sobreajuste (\textit{overfitting}) é um risco especialmente elevado em redes profundas, que possuem grande capacidade expressiva. Uma das técnicas de regularização mais eficazes para mitigá-lo é o \textit{dropout}: durante o treinamento, cada neurônio de uma camada é desativado aleatoriamente com uma probabilidade $p$, forçando a rede a aprender representações redundantes e menos dependentes de neurônios individuais \cite{bochie2020sbrc}. O \textit{dropout} é aplicado exclusivamente na fase de treinamento; na inferência, todos os neurônios são mantidos ativos e seus pesos são escalados pela probabilidade de retenção $(1 - p)$.

\section{\textit{Autoencoder}}

O Autoencoder é uma arquitetura de rede neural projetada para aprender representações compactas dos dados de forma não supervisionada \cite{bochie2020sbrc}. A rede é dividida em duas partes simétricas: o codificador (\textit{encoder}) mapeia a entrada $x$ para um vetor latente $z$ de dimensão reduzida, e o decodificador (\textit{decoder}) reconstrói a entrada $\hat{x}$ a partir de $z$. O objetivo do treinamento é minimizar o erro de reconstrução entre a entrada original e a saída reconstruída, tipicamente medido pelo erro quadrático médio:
\[
  \mathcal{L} = \frac{1}{n}\sum_{k=1}^{n}\left\|x_k - \hat{x}_k\right\|^2
\]

\textit{Autoencoders} subdimensionados possuem código de dimensão menor que a entrada, forçando o modelo a comprimir a informação e descartar ruído. \textit{Autoencoders} sobredimensionados, cujo código tem dimensão maior que a entrada, precisam de restrições adicionais para que não aprendam simplesmente a função identidade \cite{bochie2020sbrc}.

Em aplicações de detecção de anomalias, o \textit{Autoencoder} é treinado exclusivamente com amostras benignas. Durante a inferência, amostras que pertencem a classes de ataque geram erro de reconstrução elevado, pois o modelo aprendeu apenas a representar o comportamento normal. Um limiar (\textit{threshold}) é definido sobre o erro de reconstrução para separar tráfego benigno de tráfego anômalo. Essa propriedade torna o Autoencoder especialmente adequado para a detecção de ataques inéditos (\textit{zero-day}), que não possuem exemplos rotulados no conjunto de treinamento.

\section{Aplicação em Detecção de Intrusão}

Em tráfego de rede representado por vetores de características numéricas, como os derivados do CICIDS2017 e do LycoS-IDS2017 utilizados neste trabalho, a MLP classifica cada fluxo a partir de dezenas de métricas (duração, contagem de pacotes, tamanho médio, \textit{flags} TCP, entre outras). O \textit{pipeline} de treinamento envolve a coleta e rotulagem dos fluxos, o pré-processamento (normalização, tratamento de valores ausentes e balanceamento de classes), o treinamento com validação e a avaliação em um conjunto de teste retido \cite{medeiros2019sbrc}.

Duas arquiteturas são investigadas experimentalmente neste trabalho. Nos experimentos supervisionados (Baseline e Melhorias A–C), uma MLP é treinada com rótulos de classe para classificar o tráfego entre benigno e os diferentes tipos de ataque presentes nos \textit{datasets}. Nos experimentos não supervisionados (Melhorias D–G), um \textit{Autoencoder} é treinado apenas com tráfego benigno e utiliza o erro de reconstrução como pontuação de anomalia em tempo de inferência.

Uma vantagem das DNNs sobre algoritmos mais simples é a capacidade de modelar padrões complexos e não lineares no tráfego. A contrapartida é a maior demanda computacional e a menor interpretabilidade em relação ao \textit{Random Forest}. Além disso, conjuntos de dados com classes altamente desbalanceadas exigem cuidados adicionais, como reamostragem ou pesos de classe na função de perda, para que a DNN não ignore as classes de ataque minoritárias \cite{medeiros2019sbrc}. Esses aspectos são investigados nos experimentos descritos no Capítulo~7. Antes disso, o Capítulo~5 contextualiza as escolhas metodológicas à luz dos trabalhos relacionados mais relevantes da literatura, e o Capítulo~6 detalha o \textit{pipeline} experimental e os critérios de avaliação adotados.
